\chapter{بحث، نتیجه‌گیری و کارهای آینده}
\label{ch:conclusions}

\section{فاصله زیان و تولید}

یافته علمی اصلی مسیر مستقیم این است که کاهش زیان تحت تحمیل معلم از تولید
خودرگرسیو معتبر جدا است. نسخه \en{v6.2} این جداسازی را تیز می‌کند: روی همان
\N{32} نمونه آموزشی، زیان متنی \N{32.20} درصد کاهش یافت اما رفتار آموخته‌شده
به تولید قابل اتکا منتقل نشد. توضیح سازگار با شواهد، اریبی مواجهه است نه
نبود ظرفیت. پیامد روشمند آن است که انتخاب چک‌پوینت گفتار‌به‌گفتار نباید تنها
بر زیان آفلاین استوار باشد.

در مسیر آبشاری، دو جداسازی مشابه دیده شد. نخست، پنل مکانیکی \N{9/9} با
پاسخ‌گوی \N{0.5B} کیفیت معنایی پاسخ‌گوی \N{4B} را ثابت نمی‌کند. دوم، عبور
آزمون خودکار هم‌خانواده کیفیت ادراکی را ثابت نمی‌کند. فاصله میان نرخ خطای
نویسه \N{7.14} درصد و نرخ خطای واژه \N{30.41} درصد نیز نشان می‌دهد که حفظ
تقریبی نویسه‌ها با فهم‌پذیری واژه‌ای یکی نیست. امتیاز پایین مرجع هم‌ترازشده در آزمون معنایی یادآوری می‌کند که نوبت بعدی
پیکره لزوماً پاسخ مناسب به آخرین اظهار کاربر نیست؛ به همین دلیل آداپتر
\en{LoRA} که زیان توسعه را کاهش داد، از خط پایه پیشی نگرفت.

\section{علت شکست مسیر مستقیم}

ممیزی پایانی نشان داد که توکن‌ساز انگلیسی \en{Moshi} برای هر واژه فارسی \N{6.6}
نشانه می‌خواهد و \N{98.5} درصد نوبت‌ها از بودجه یک نشانه در قاب فراتر می‌روند.
برداشتن این گلوگاه با لاتین‌نویسی واجی بهترین نتیجه تولید را از \N{3/9} به
\N{6/9} رساند و خروجی ساکت را شنیدنی کرد، اما گفتار فهم‌پذیر نشد؛ پس گلوگاه
توکن‌ساز واقعی است اما تنها مانع نیست. \en{J-Moshi}~\cite{ohashi2025jmoshi} نیز
برای انتقال \en{Moshi} به ژاپنی توکن‌ساز را جایگزین کرد و سپس به حدود
\N{60,000} ساعت پیش‌آموزش گفت‌وگو نیاز داشت، بیش از پانصد برابر داده این پروژه. مدل ورودی گفتاری \en{Gemma~4} نیز نشان
داد که حذف بازشناسی، حتی با همان مدل زبانی، کیفیت را کم می‌کند.

\section{محدودیت‌ها و ملاحظات}

شواهد تأخیر و وقفه خودکار و روی \en{H100} است: حافظه یک کارت \N{12}
گیگابایتی رعایت شد اما سرعت آن نه، و کاربر انسانی، مسیر فیزیکی میکروفون و
بلندگو و پژواک واقعی در کار نبود. برچسب‌های وقفه و داده خودکارند و برچسب
انسانی ندارند؛ دو داور خودکار پاسخ‌گو اختلاف زیادی دارند؛ صدای
سالمند در آزمون‌ها شبیه‌سازی است. تشخیص وقفه به بازشناس گفتار تکیه دارد، نه
به ویژگی‌های صوتی تعیین‌شده در تعریف پروژه. مجموعه داده از آرشیو داخلی مشتق
شده و قابل بازنشر نیست. نتایج مطالعه شنیداری هنوز دریافت نشده است.

\section{پاسخ به پرسش‌های پژوهش}

\مهم{۱.} تطبیق کم‌پارامتر \en{Moshika} سرگشتگی آزمون کنارگذاشته را کاهش داد
(نشانه متنی \N{2.53} به \N{1.19}، صوتی \N{13.98} به \N{4.71}) اما گفتار
فهم‌پذیر نساخت؛ توکن‌ساز انگلیسی مانع اصلی است. \مهم{۲.} در مسیر آبشاری،
تشخیص پیش‌دستانه پایان گفتار با نگه‌داشت صوت تا تأیید، تولید جریانی با قید
خط فارسی، بندبندی، صف پخش و کنترل وقفه دومرحله‌ای با تاریخچه شنیده‌شده، تأخیر را
به میانگین \N{372} و صدک نود \N{431} میلی‌ثانیه رساند. \مهم{۳.} اندازه کیفیت
نخرید: کوچک‌ترین بازشناس‌ها و گفتارساز تخصصی فارسی بهترین بده‌بستان را داشتند،
پاسخ‌گوهای \N{9} تا \N{36} میلیاردی زیر هر دو داور از بهترین مدل چهارمیلیاردی
بهتر نبودند، و مدل سرتاسری \N{7.7} میلیارد پارامتری به گفتار فهم‌پذیر نرسید. \مهم{۴.} درک
سالمندان به نتایج مطالعه شنیداری وابسته است. وضعیت هر نیازمندی در
\mbox{جدول~\ref{tab:req-status}} آمده است.

\section{کارهای آینده}

تکمیل شواهد: اجرای همان آزمون حلقه‌ای روی کارت \N{12} تا \N{24} گیگابایتی و
با مسیر فیزیکی صدا؛ گردآوری پاسخ‌های مطالعه شنیداری و یک جلسه تعامل زنده با
سالمندان؛ بازبینی انسانی نمونه‌ای از برچسب‌ها؛ و داوری انسانی پاسخ‌گوها
(ترجیح شرکت‌کنندگان درباره جایگزینی \en{Qwen3.5-4B} تصمیم می‌گیرد). بهبود سامانه:
جداکردن آموخته واژه شنونده از محتوایی در آغاز بریده‌شده (پیش‌بافر \N{200}
میلی‌ثانیه‌ای آزموده و رد شد)، تشخیص پایان نوبت معنایی برای درنگ‌های گفتار سالمندان
(در دو دور آزموده شد و نوبت‌های تمام‌شده بسیاری را به تأخیر انداخت؛ نگه‌داشتن
کوتاه‌تر یا برچسب انسانی پایان نوبت لازم است)، و بازشناس
دقیق‌تر در همان بودجه تأخیر. مسیر مستقیم: توکن‌سازی متناسب با بودجه قاب و
راهکاری برای شکاف واداشت معلم و تولید آزاد، و رویکرد مدل منجمد.
